\subsection{Weyl 群}

设 G 是连通紧李群，T 是 G 的极大环面。用 \(N_{G}(T)\) 表示 T 在 G 中的正规化子；由第 4 节命题 4，商群 \(N_{G}(T)/T\) 有限。我们把它记作 \(W_{G}(T)\) 或 \(W(T)\) ，称之为 G 的极大环面 T 的 Weyl 群，或 G 相对于 T 的 Weyl 群。由于 T 交换，\(N_{G}(T)\) 经 G 的内自同构在 T 上的作用过渡到商，便得到群 \(W_{G}(T)\) 在李群 T 上的一个作用，称为典则作用。由第 2 节定理 2 的推论 6，这个作用是忠实的：相伴的同态 \(W_{G}(T) \rightarrow Aut T\) 是单射。

若 \(T'\) 是 G 的另一个极大环面，且 \(g \in G\) 使得 Int g 把 T 映为 \(T'\) （第 2 节，定理 2 b)），则 Int g 诱导一个同构 \(a_g\) （从 \(\mathrm{W}_{\mathrm{G}}(\mathrm{T})\) 到 \(\mathrm{W}_{\mathrm{G}}(\mathrm{T}')\) ），并且 \(a_g(s)(gtg^{-1}) = gs(t)g^{-1}\) 对一切 \(s \in \mathrm{W}_{\mathrm{G}}(\mathrm{T})\) 与一切 \(t \in T\) 成立。

命题 5. a) G 的每个共轭类都与 T 相交。

\begin{enumerate}
\def\labelenumi{\alph{enumi})}
\setcounter{enumi}{1}
\tightlist
\item
  \(\mathrm{T}\) 与 \(\mathbf{G}\) 的诸共轭类的交是 Weyl 群的轨道。
\end{enumerate}

设 \(g \in \mathbf{G}\) ；由第 2 节定理 2，存在 \(h \in \mathbf{G}\) 使得 \(g \in h\mathrm{Th}^{-1}\) ，由此得 \(a\) )。由 Weyl 群的定义，\(\mathrm{W_G(T)}\) 在 \(T\) 上同一轨道中的任两个元素在 \(G\) 中共轭；反之，设 \(a, b\) 是 \(T\) 的两个在 \(G\) 中共轭的元素。存在 \(h \in G\) 使得 \(b = hah^{-1}\) ；对 \(A = \{a\}\) 、\(s = \operatorname{Int} h\) 、\(T' = T\) 应用第 2 节定理 2 的推论 7，可见存在 \(g \in G\) 使得 \(\operatorname{Int} hg\) 把 \(T\) 映到 \(T\) 且把 \(a\) 映到 \(b\) 。于是 \(hg\) 在 \(\mathrm{W_G(T)}\) 中的类把 \(a\) 映到 \(b\) ，命题得证。

推论 1. T 到 G 中的典则单射经过渡到商，定义一个从 \(\mathrm{T}/\mathrm{W}_{\mathrm{G}}(\mathrm{T})\) 到 G 的共轭类空间 \(\mathrm{G}/\mathrm{Int}(\mathrm{G})\) 上的同胚。

事实上，这是两个紧空间之间的连续双射（参见《一般拓扑学》，第 III 章，第 29 页，推论 1）。

推论 2. 设 E 是 G 的在内自同构下稳定的子集。则 E 在 G 中开（相应地，闭、稠密）当且仅当 \(E \cap T\) 在 T 中开（相应地，闭、稠密）。

这由推论 1 以及典则映射 \(T \rightarrow T/W_{G}(T)\) 与 \(G \rightarrow G/Int(G)\) 是开映射得出（《一般拓扑学》，第 III 章，第 10 页，引理 2）。

用 g 表示 G 的李代数，用 t 表示 T 的李代数。\(W_{\mathrm{G}}(\mathrm{T})\) 在 T 上的作用诱导出群 \(W_{\mathrm{G}}(\mathrm{T})\) 在 R-向量空间 t 上的一个表示，称为典则表示。

命题 6. a) G 在 g 上（对伴随表示而言）的每条轨道都与 t 相交。

\begin{enumerate}
\def\labelenumi{\alph{enumi})}
\setcounter{enumi}{1}
\tightlist
\item
  \(\mathfrak{t}\) 与 \(\mathrm{G}\) 的诸轨道的交是 \(\mathrm{W_G(T)}\) 在 \(\mathfrak{t}\) 上的轨道。
\end{enumerate}

断言 \(a\) ) 由第 1 节定理 1 得出。设 \(x, y\) 是 \(\mathfrak{t}\) 的两个在 \(\operatorname{Ad}(\mathrm{G})\) 下共轭的元素，\(h \in \mathrm{G}\) 满足 \((\operatorname{Ad} h)(x) = y\) 。对 \(\mathfrak{a} = \{x\}\) 、\(u = \operatorname{Ad} h, \mathfrak{t}' = \mathfrak{t}\) 应用第 1 节定理 1 的推论，可见存在 \(g \in \mathrm{G}\) 使得 \(\operatorname{Ad} hg\) 把 \(\mathfrak{t}\) 映到 \(\mathfrak{t}\) 且把 \(x\) 映到 \(y\) 。于是 \(hg \in \mathrm{N}_{\mathrm{G}}(\mathrm{T})\) （第 III 章，§9，第 4 节，命题 11），且 \(hg\) 在 \(\mathrm{W}_{\mathrm{G}}(\mathrm{T})\) 中的类把 \(x\) 映到 \(y\) ，命题得证。

推论. t 到 g 中的典则单射经过渡到商，定义一个从 \(\mathfrak{t}/\mathrm{W}_{\mathrm{G}}(\mathrm{T})\) 到 \(\mathfrak{g}/\mathrm{Ad}(\mathrm{G})\) 上的同胚。

把这个映射记作 \(j\) ；它是双射且连续（命题 6）。我们有交换图式

\[
\begin{array}{c c c} \mathfrak {t} & \stackrel {{i}} {{\longrightarrow}} & \mathfrak {g} \\ p \Big \downarrow & & \Big \downarrow q \\ \mathfrak {t} / \mathrm{W} _ {\mathrm{G}} (\mathrm{T}) & \stackrel {{j}} {{\longrightarrow}} & \mathfrak {g} / \mathrm{Ad} (\mathrm{G}) \end{array}
\]

其中 p 与 q 是商映射，i 是典则单射。由于 i 与 q 是真的（《一般拓扑学》，第 I 章，§10，第 1 节，命题 2，及《一般拓扑学》，第 III 章，§4，第 1 节，命题 2 c)），且 p 是满射，由此得出 j 是真的（《一般拓扑学》，第 I 章，§10，第 1 节，命题 5），从而是同胚。

命题 7. 设 H 是 G 的包含 T 的闭子群。

\begin{enumerate}
\def\labelenumi{\alph{enumi})}
\item
  用 \(\mathrm{W}_{\mathrm{H}}(\mathrm{T})\) 表示子群 \(\mathrm{N}_{\mathrm{H}}(\mathrm{T})/\mathrm{T}\) （\(\mathrm{W}_{\mathrm{G}}(\mathrm{T})\) 中的）；群 \(H/H_{0}\) 同构于商群 \(\mathrm{W}_{\mathrm{H}}(\mathrm{T})/\mathrm{W}_{\mathrm{H}_{0}}(\mathrm{T})\) 。
\item
  H 连通当且仅当 \(\mathrm{W}_{\mathrm{G}}(\mathrm{T})\) 中在 H 有代表的每个元素都属于 \(\mathrm{W}_{\mathrm{H}_{0}}(\mathrm{T})\) 。
\end{enumerate}

断言 \(a\) ) 由第 2 节定理 2 的推论 8 得出，而断言 \(b\) ) 是 \(a\) ) 的一个特殊情形。

